\section{Reward Function}
\label{sec:reward}

\subsection{Design philosophy}
Early experiments using a single $\{0,1\}$ correctness reward produced
unstable training: gradient signals were sparse and the policy quickly
learnt to exploit boundary cases (truncated answers, role-token leakage,
or extremely long output that hides the answer). Following the lesson
patterns in DeepSeek-Math~\cite{deepseekmath} and DAPO~\cite{dapo}, we use a
\emph{composite} reward that mixes positive shaping terms for correct and
clean outputs with several explicit penalties.

The total reward $R(q,o,a^{*})$ for a prompt $q$, completion $o$, and
reference answer $a^{*}$ is the sum of all components defined in
\cref{tab:reward-components}.

\subsection{Components}
\begin{table}[h]
\centering
\small
\begin{tabularx}{\linewidth}{lrX}
\toprule
\textbf{Component} & \textbf{Score} & \textbf{When applied} \\
\midrule
$\rExact$       & $+1.00$         & extracted final answer is character-equal to $a^{*}$ after normalisation \\
$\rEquiv$       & $+0.80$         & extracted answer is mathematically equivalent to $a^{*}$ but textually different \\
$\rFmt$         & $+0.05$         & output ends with a single well-formed \code{Final Answer:} line \\
$\rClean$       & $+0.03$         & no spurious trailing content after the answer line \\
\code{<answer>} bonus & $0.00$    & reserved hook, currently disabled \\
$\pNoAns$       & $-0.12$         & no answer can be extracted from the output \\
$\pRep$         & $-0.08$         & visible repetition (n-gram or paragraph) detected \\
$\pRole$        & $-0.08$         & role tokens (\code{user}, \code{assistant}, \dots) leak into the output \\
$\pTail$        & up to $-0.20$   & long irrelevant text after the answer line \\
$\pLen$         & up to $-0.05$   & output exceeds the length budget \\
\bottomrule
\end{tabularx}
\caption{All reward components used in training. All shape terms are
additive.}
\label{tab:reward-components}
\end{table}

\subsection{Formal expression}
Let $\hat a(o)$ denote the answer extracted from the completion $o$. With
the indicator $\mathbb{1}[\cdot]$ we can write
\begin{align}
R(q, o, a^{*})
&= \rExact \cdot \mathbb{1}[\hat a(o) \equiv_{\text{str}} a^{*}]
 + \rEquiv \cdot \mathbb{1}[\hat a(o) \equiv_{\text{math}} a^{*}] \notag\\
&\;{+}\; \rFmt \cdot \mathbb{1}[\text{format ok}]
 + \rClean \cdot \mathbb{1}[\text{no tail}] \notag\\
&\;{-}\; \pNoAns \cdot \mathbb{1}[\hat a(o) = \varnothing]
 - \pRep   \cdot \mathbb{1}[\text{repetition}] \notag\\
&\;{-}\; \pRole - \pTail(o) - \pLen(o),
\label{eq:reward}
\end{align}
where the two correctness indicators are mutually exclusive (an exact
match takes precedence over equivalence). The length penalty $\pLen(o)$
starts at zero up to a budget of \textbf{256 tokens} and grows linearly
beyond that, capped at $-0.05$.

\subsection{Answer extraction and normalisation}
Three answer formats are accepted in parallel:
\code{Final Answer: <answer>}, \code{<answer>...</answer>}, and
\code{$\backslash$boxed\{...\}}. After extraction the answer string goes
through a normalisation pass that
(i)~removes thousand separators in numbers,
(ii)~unifies $\sqrt{\cdot}$ written as \code{\textbackslash sqrt\{\}} or as
the Unicode glyph \code{$\sqrt{\,}$},
(iii)~normalises common minus and division signs,
(iv)~accepts simple algebraic equivalents (\eg $0.5$ vs.\ $1/2$), and
(v)~trims trailing text past the answer anchor when present. Equivalence
matching is checked after normalisation only when the exact match fails.

\subsection{Empirical observations}
The composite design has two practical effects. First, removing $\pTail$
in an ablation produced visible reward hacking: long noise after a correct
answer line. Second, removing $\pLen$ slowly inflated average response
length without a measurable accuracy gain, wasting decoding budget. Both
penalties stayed in the final recipe.
